% ============================================================
% DGF Supplemental Sections: S1 Reflux Bound + S4 BH Entropy
% PRD revtex4-2 format.  Insert into main PRD manuscript or
% compile standalone as Supplemental Material.
% ============================================================
\documentclass[aps,prd,twocolumn,superscriptaddress,nofootinbib,floatfix]{revtex4-2}
\usepackage{amsmath,amssymb,amsthm}
\usepackage{graphicx}
\usepackage[colorlinks=true,citecolor=blue,linkcolor=blue,urlcolor=blue]{hyperref}

% Theorem environments
\newtheorem{theorem}{Theorem}[section]
\newtheorem{lemma}[theorem]{Lemma}
\newtheorem{corollary}[theorem]{Corollary}
\newtheorem{conjecture}[theorem]{Conjecture}

\begin{document}

\title{Supplemental Material: \\
Reflux Bound and Black-Hole Entropy Area Law \\
in the Decoherence Geometry Framework}
\author{Huang Zhongchang}
\affiliation{Independent Researcher, Nanjing, China}
\email{valhuang@kaiwucl.com}
\date{\today}

\maketitle

% ============================================================
\section{Reflux Bound: Capacity Limit on Information Backflow}
\label{sec:S1}
% ============================================================

\subsection{Physical setup}

Consider a system $S$ partitioned into $N_S$ Planck cells and an environment
$E$ with $N_E$ cells, each of volume $\ell^3$.  Within the DGF binary pointer
basis, every cell $c$ is a qubit:
\begin{equation}
  |0\rangle_c = \text{flowing (accessible) coherence},\qquad
  |1\rangle_c = \text{archived (locked) information}.
  \label{eq:S1_basis}
\end{equation}
The accessible coherence fraction of any subsystem $X$ is the ensemble average
of cell occupancies,
\begin{equation}
  q_X = \frac{1}{N_X}\sum_{c\in X} \langle 0|\rho_c|0\rangle
      = 1 - A_X,
  \label{eq:S1_qX}
\end{equation}
with $\rho_c$ the reduced state of cell $c$ and $A_X$ the archived fraction.
Equation~\eqref{eq:S1_qX} is the operational DGF definition~\cite{dgf_prd}
and requires no spacetime or mass input.

\subsection{GKSL directionality postulate}

The dynamics of the joint $S+E$ system is generated by a Gorini--Kossakowski--
Sudarshan--Lindblad (GKSL) master equation.  The central Tier-2 postulate is
that the Lindblad generator $L$ is \emph{directed} with respect to the
$\{|0\rangle,|1\rangle\}$ basis:
\begin{equation}
  \boxed{
  L_k = \sqrt{\gamma_k}\,|1\rangle\langle 0|,\qquad \gamma_k \ge 0,
  }
  \label{eq:S1_L_directed}
\end{equation}
for every Lindblad operator $L_k$.  No operator maps $|1\rangle\to|0\rangle$
at the bare-GKSL level.  Equation~\eqref{eq:S1_L_directed} guarantees that the
ensemble-mean flow of information is from accessible ($|0\rangle$) to archived
($|1\rangle$) --- decoherence is irreversible in the thermodynamic limit.

\subsection{Capacity definitions}

\begin{definition}[Forward capacity]
The forward information-transport capacity of the environment is the total
number of accessible ($|0\rangle$) qubits it can receive:
\begin{equation}
  C_F = N_E\,q_E.
  \label{eq:S1_CF}
\end{equation}
\end{definition}

\begin{definition}[Backflow capacity constraint]
Any reflux event requires a $|1\rangle$ in $E$ to transition back to
$|0\rangle$.  The receiving capacity on the system side is at most its total
accessible count:
\begin{equation}
  N_{\rm back} \le N_S\,q_S.
  \label{eq:S1_Nback}
\end{equation}
\end{definition}

\begin{theorem}[Reflux Bound]
\label{thm:reflux}
For $N_S = N_E$, the probability $P_{\rm reflux}$ that any single
information-transfer event constitutes a backflow from $E$ to $S$ satisfies
\begin{equation}
  \boxed{
  P_{\rm reflux} \le \frac{q_S}{q_E}.
  }
  \label{eq:S1_reflux_bound}
\end{equation}
\end{theorem}

\begin{proof}
Forward flow and backflow compete for the same information channels.
In the stationary ensemble the rate of forward events is proportional to
$C_F = N_E q_E$ while the rate of backflow events is bounded by
$N_S q_S$.  Hence
\begin{equation}
  P_{\rm reflux}
  = \frac{\text{backflow rate}}{\text{total transfer rate}}
  \le \frac{N_S q_S}{N_E q_E}
  = \frac{q_S}{q_E},
  \label{eq:S1_proof}
\end{equation}
where the last equality uses $N_S=N_E$.
\end{proof}

\subsection{Domain of applicability}

The bound is meaningful only when:
\begin{enumerate}
  \item $q_E > 0$ --- the environment retains some accessible coherence;
    otherwise forward transport cannot be defined.
  \item Irreversible forward transport \emph{has occurred}, so that $C_F$ is
    populated and $P_{\rm reflux}$ is a well-posed conditional probability.
  \item $q_S < q_E$ --- when the system is \emph{more classical} than the
    environment, the bound is $<1$ and nontrivial.  For $q_S \ge q_E$ the
    bound reduces to $P_{\rm reflux}\le 1$, which is always satisfied and
    carries no physical content.
\end{enumerate}

\subsection{Relation to BLP detection and Buscemi classification}

The Reflux Bound complements two established frameworks for quantifying
non-Markovianity:

\begin{itemize}
  \item \textbf{BLP measure} (Breuer--Laine--Piilo~\cite{blp2009}):
    detects reflux by monitoring whether the trace distance
    $D[\rho_1(t),\rho_2(t)]$ increases in time.  BLP answers
    \emph{whether} backflow occurs.  Theorem~\ref{thm:reflux} answers
    \emph{how much} backflow is possible, purely from the accessible-coherence
    inventory of sender and receiver --- no dynamical integration required.
  \item \textbf{Buscemi classification}~\cite{buscemi2014}:
    decomposes information flow into entanglement, quantum discord, and
    classical correlation channels.  Theorem~\ref{thm:reflux} bounds the
    \emph{total magnitude} of any such classified backflow channel.  When
    $q_S/q_E \ll 1$, no Buscemi-classified channel can carry substantial
    reflux, regardless of its internal quantum or classical character.
\end{itemize}

In short: BLP provides the detector, Buscemi provides the taxonomy, and
S1 provides the capacity-theoretic magnitude ceiling.

\subsection{$T_1$ correction}

Equation~\eqref{eq:S1_L_directed} is exact only in the limit of infinite
relaxation time $T_1\to\infty$.  For any finite $T_1$, spontaneous
$|1\rangle\to|0\rangle$ relaxation occurs with characteristic timescale
$T_1$, introducing a correction
\begin{equation}
  \delta \sim \frac{\tau}{T_1},
  \label{eq:S1_delta}
\end{equation}
where $\tau$ is the GKSL coarse-graining timescale.  The refined bound reads
\begin{equation}
  P_{\rm reflux} \le \frac{q_S}{q_E} + O\!\left(\frac{\tau}{T_1}\right).
  \label{eq:S1_refined}
\end{equation}
In macroscopic environments ($T_1$ effectively infinite on laboratory
scales) the correction is negligible.

\subsection{Generator decomposition}

The Lindblad generator admits the orthogonal decomposition
\begin{equation}
  L = L_{\rm copy} \oplus L_{\rm swap},
  \label{eq:S1_decomp}
\end{equation}
where $L_{\rm copy}$ contains only $|0\rangle\!\to\!|1\rangle$ operators
(purely directional, implementing Assumption~A3 at the dynamical level) and
$L_{\rm swap}$ contains Hermitian-preserving operators that exchange
$|0\rangle\leftrightarrow|1\rangle$ (enabling backflow).  Theorem~\ref{thm:reflux}
constrains the \emph{relative strength} of the swap sector: writing
$\|L_{\rm swap}\|/\|L_{\rm copy}\| = \varepsilon$, the backflow probability
scales as $\varepsilon^2$, and the capacity argument requires
$\varepsilon^2 \le q_S/q_E$.

\subsection{Assumption tiering (Tier 0--3)}

\begin{description}
  \item[Tier 0 --- Founding axioms:]
    A1 (information exists), A2 (capacity bounded at scale $\ell$),
    A3 (overflow irreversible).  No spacetime, no mass, no $G$.

  \item[Tier 1 --- Theorems from Tier 0 alone:]
    Definition of $q$ (Eq.~\ref{eq:S1_qX}), capacity formulas
    $C_F=N_E q_E$ and $N_{\rm back}\le N_S q_S$ (Eqs.~\ref{eq:S1_CF},
    \ref{eq:S1_Nback}).  The inequality $P_{\rm reflux}\le q_S/q_E$
    follows from Tier~0~+~1 without additional physical postulates.

  \item[Tier 2 --- Additional physical postulates:]
    \begin{itemize}
      \item \textbf{Qubit binary pointer basis} (Eq.~\ref{eq:S1_basis}):
        each Planck cell is a two-level system.  This is \emph{not} derived
        from A1--A3; it is a modeling postulate consistent with the capacity
        bound A2 ($O(1)$ bits per cell).
      \item \textbf{GKSL directionality} (Eq.~\ref{eq:S1_L_directed}):
        $L$ only creates $|1\rangle$ from $|0\rangle$.  This is a dynamical
        implementation of A3; it is \emph{not} forced by A3 alone, since A3
        is a thermodynamic statement while~\eqref{eq:S1_L_directed} is a
        microscopic operator constraint.
      \item \textbf{Equal cell count} $N_S=N_E$: a normalization condition
        for the comparison; the general bound is $P_{\rm reflux}\le
        (N_S q_S)/(N_E q_E)$.
    \end{itemize}

  \item[Tier 3 --- Modeling choices and idealizations:]
    \begin{itemize}
      \item \textbf{1-bit per cell}: A2 states ``at most $O(1)$ bits''; the
        qubit model fixes this to exactly one bit.  This is a \emph{postulate},
        not a derived result.
      \item \textbf{Pure-dephasing $H_{\rm int}$}: the interaction Hamiltonian
        is taken to commute with the pointer observable, so that the GKSL
        directionality~\eqref{eq:S1_L_directed} is not undermined by unitary
        $|1\rangle\!\to\!|0\rangle$ rotations.  Real interactions contain
        small non-commuting components; their omission is an idealization.
      \item \textbf{Markovian bath}: the GKSL form assumes
        $\tau_{\rm corr}^{\rm(E)} \ll \tau_S$, i.e., environment correlations
        decay faster than any system timescale.
    \end{itemize}
\end{description}

\subsection{Honest summary}

\begin{enumerate}
  \item The inequality $P_{\rm reflux}\le q_S/q_E$ follows from
    capacity counting (Tier~0+1).  It is \emph{not} a dynamical derivation
    from a microscopic Hamiltonian.
  \item The qubit pointer basis (1-bit postulate) is an \textbf{assumption},
    not a theorem.  The bound is proportionally weaker if cells carry $b>1$
    bits and only one bit participates in the pointer.
  \item The GKSL directionality~\eqref{eq:S1_L_directed} is a
    \textbf{modeling choice} that enforces A3 at the operator level.  A3
    alone (thermodynamic irreversibility) does not logically entail
    Eq.~\eqref{eq:S1_L_directed}.
  \item Pure-dephasing $H_{\rm int}$ is an \textbf{idealization}.  Real
    systems have small non-commuting interaction terms that can induce
    unitary back-rotation; these are captured only through the $O(\tau/T_1)$
    correction.
  \item The bound is \textbf{nontrivial only when $q_S<q_E$}.  For $q_S\ge q_E$
    it reduces to $P_{\rm reflux}\le 1$, which is tautological.
  \item The result should be read as: \emph{If} decoherence is governed by a
    directed binary pointer basis with GS KL dynamics, \emph{then} the
    probability of reflux cannot exceed the accessible-coherence ratio.
    This is a conditional theorem, not an unconditional law of nature.
\end{enumerate}

% ============================================================
\section{Black-Hole Entropy: Min-Cut Area Law from DGF}
\label{sec:S4}
% ============================================================

\subsection{Core result}

Within the DGF information graph, the entropy of a bounded region $\Omega$
as seen from its exterior is bounded by the number of Planck cells on the
boundary $\partial\Omega$:
\begin{equation}
  \boxed{
  S_{\rm DGF}(\Omega) \le N_{\partial\Omega}\,\ln 2
  \propto \frac{A}{\ell^2}.
  }
  \label{eq:S4_main}
\end{equation}
This is the Bekenstein--Hawking area law, obtained from graph theory
(Ford--Fulkerson min-cut) plus the two founding DGF axioms A1 (information
exists) and A2 (capacity is bounded).  No gravitational field equations,
no event-horizon thermodynamics, and no spacetime metric are used in the
derivation of the inequality.  The proportionality $N_{\partial\Omega}\propto A$
is a theorem; the numerical coefficient is model-dependent and discussed
in~\S\ref{sec:S4_coefficient}.

\subsection{Layer 1: Min-cut bound (strict graph theory)}

\label{sec:S4_layer1}

Let $G=(V,E)$ be the DGF information graph whose vertices are Planck cells
$c\in V$ and whose edges connect cells that are adjacent in the
coarse-grained continuum limit.  A \emph{cut} $\partial\Omega$ is a set of
edges whose removal disconnects the interior $\Omega$ from the exterior
$V\setminus\Omega$.

\begin{lemma}[Ford--Fulkerson~\cite{ford1956}]
Let $C(\partial\Omega)$ be the capacity of cut $\partial\Omega$, defined as
the number of edges crossing it (each edge carries $\le 1$ bit by A2).
Then the maximum information flow from $\Omega$ to $V\setminus\Omega$ equals
the capacity of the minimum cut:
\begin{equation}
  \max\text{-flow}(\Omega\to V\setminus\Omega)
  = \min_{\partial\Omega'}\,C(\partial\Omega').
  \label{eq:S4_ff}
\end{equation}
\end{lemma}

\begin{theorem}[Min-Cut Entropy Bound]
\label{thm:mincut}
The number of internal configurations of $\Omega$ that are distinguishable
from the exterior is bounded by $2^{|\partial\Omega|}$, where
$|\partial\Omega| = N_{\partial\Omega}$ is the number of edges on the
minimum cut.  Hence
\begin{equation}
  S_{\rm DGF}(\Omega) \le N_{\partial\Omega}\,\ln 2.
  \label{eq:S4_mincut}
\end{equation}
\end{theorem}

\begin{proof}
Each edge on the minimum cut can transmit at most one bit (A2).  The
exterior can therefore discriminate at most $2^{N_{\partial\Omega}}$
distinct internal states.  By Boltzmann's entropy formula
$S = \ln(\text{\# accessible microstates})$, the bound follows.
No property of the interior dynamics is used --- only the graph topology
and the capacity bound A2.  The inequality~\eqref{eq:S4_mincut} is
therefore a \textbf{strict graph-theoretic theorem} from A1~+~A2 alone.
\end{proof}

The boundary cell count scales with area:
\begin{equation}
  N_{\partial\Omega} = \frac{|\partial\Omega|}{\ell^2}
  = \frac{A}{\ell^2},
  \label{eq:S4_Narea}
\end{equation}
where $A$ is the area of the bounding surface in the continuum limit.
Thus $S_{\rm DGF} \propto A$ --- the area law --- is obtained without using
any property of gravity.  It is a consequence of bounded information capacity
on a graph.

\subsection{Layer 2: $q$-field determines the cut surface}
\label{sec:S4_layer2}

Layer~1 establishes $S\le N_{\partial\Omega}\propto A$ for \emph{any} cut.
Which cut corresponds to the physical black-hole horizon is determined by the
$q$-field.

The static, spherically symmetric $q$-field solution for a point mass $M$ is
(from the DGF entropy-variational principle in the $q\ll 1$ regime):
\begin{equation}
  q(r) = q_\infty - \frac{M}{4\pi m_0}\,\frac{\ell}{r},
  \label{eq:S4_qfield}
\end{equation}
with $m_0 = \hbar/(c\ell)$ the derived mass scale.  The information horizon
forms where coherent transmission across the surface becomes exponentially
suppressed.  A natural criterion is that the transmission amplitude
$\propto q$ drops to $e^{-1/2}\approx 0.607$, i.e.,
\begin{equation}
  q(r_s) = e^{-1/2}\quad\Longrightarrow\quad
  r_s = \frac{M\ell}{4\pi m_0}\,
        \frac{1}{q_\infty - e^{-1/2}}.
  \label{eq:S4_rs}
\end{equation}

This is the DGF prediction for the Schwarzschild radius.  Crucially, the
$q$-field \emph{locates} the cut --- it tells us \emph{where} the min-cut
sits --- but the entropy bound~\eqref{eq:S4_mincut} does not depend on the
precise value of $r_s$; it only requires that a cut exists and that its area
scales as $r_s^2$.  The location determines the coefficient when comparing
with Bekenstein--Hawking.

\subsection{Layer 3: Coefficient comparison with Bekenstein--Hawking}
\label{sec:S4_coefficient}

The Bekenstein--Hawking entropy is
\begin{equation}
  S_{\rm BH} = \frac{k_B c^3 A}{4G\hbar}
             = \frac{A}{4\ell_P^2},
  \label{eq:S4_SBH}
\end{equation}
where $\ell_P = \sqrt{\hbar G/c^3}$ and we work in units $k_B=1$.
The DGF entropy bound, with cutoff scale $a$ (identifying each boundary cell
with area $a^2$ on the cut surface), is
\begin{equation}
  S_{\rm DGF} = N_{\partial\Omega}\,\ln 2
              = \frac{A}{a^2}\,\ln 2.
  \label{eq:S4_SDGF}
\end{equation}
Their ratio is
\begin{equation}
  \frac{S_{\rm DGF}}{S_{\rm BH}}
  = 4\ln 2\left(\frac{\ell_P}{a}\right)^{\!2}.
  \label{eq:S4_ratio}
\end{equation}

Two calibrations of the cutoff $a$ are physically relevant:

\begin{enumerate}
  \item \textbf{$a = \ell_P$ (Planck cutoff).}
    The DGF graph spacing is identified with the Planck length itself.
    Then
    \begin{equation}
      \frac{S_{\rm DGF}}{S_{\rm BH}} = 4\ln 2 \approx 2.77.
      \label{eq:S4_ratio1}
    \end{equation}
    The DGF bound is a factor $\sim 2.8$ above Bekenstein--Hawking.
    This is a non-trivial consistency check: the information-theoretic
    bound is in the same order of magnitude as the gravitational entropy,
    despite being derived without $G$.

  \item \textbf{$S_{\rm DGF}=S_{\rm BH}$ (matching condition).}
    Setting the ratio to unity fixes
    \begin{equation}
      a = \sqrt{4\ln 2}\; \ell_P \approx 1.66\,\ell_P.
      \label{eq:S4_matching}
    \end{equation}
    This is within a factor $<2$ of the Planck scale.  If the DGF
    fundamental scale $\ell$ differs from $a$ by an $O(1)$ geometric
    factor (e.g., from the specific lattice realization of the continuum
    limit), the Bekenstein--Hawking coefficient is recovered without
    tuning.
\end{enumerate}

The ratio $S_{\rm DGF}/S_{\rm BH}$ is therefore \textbf{model-dependent}
at the $O(1)$ level.  It depends on the cutoff scale $a$, which is not
independently predicted by the current version of DGF (see Gap~6 of the
main text for the continuum-limit universality class).  The \emph{scaling}
$S\propto A$ is a theorem; the \emph{coefficient} is a consistency
condition that constrains the cutoff.

\subsection{Layer 4: Causal-locking saturation conjecture (C1)}
\label{sec:S4_layer4}

The inequality~\eqref{eq:S4_mincut} is an upper bound.  Whether the bound is
saturated --- i.e., whether $S_{\rm DGF}(\Omega) = N_{\partial\Omega}\ln 2$
--- requires additional physics.

\begin{conjecture}[Causal Locking Saturation, C1]
\label{conj:C1}
For a region $\Omega$ whose $q$-field satisfies $q(r_s)\le e^{-1/2}$ on
its boundary, all interior lattice degrees of freedom that are not fixed
by the boundary condition are \emph{causally suppressed}: the
$q$-field's directionality (A3) prevents independent interior
fluctuations from propagating to the exterior.  Consequently, every
boundary edge carries its full A2 capacity of one bit, and the min-cut
bound is saturated:
\begin{equation}
  S_{\rm DGF}(\Omega) = N_{\partial\Omega}\,\ln 2
  \qquad\text{(if C1 holds)}.
  \label{eq:S4_saturated}
\end{equation}
\end{conjecture}

\textbf{Status of C1.}  C1 is \textbf{not proven}.  It is a physically
motivated conjecture based on the following reasoning: (i)~A3 ensures
overflow is unidirectional; (ii)~when $q(r_s)\le e^{-1/2}$, coherent
transmission across the horizon is suppressed by a factor $e^{-1}$;
(iii)~interior cells that cannot signal to the exterior cannot contribute
independent distinguishability, so the counting saturates at the boundary.
However:
\begin{itemize}
  \item The $e^{-1/2}$ threshold is a \emph{choice}, not a derived critical
    value.  It is motivated by the $1/e$ transmission-power criterion but
    other thresholds (e.g., $q=1/2$, the percolation threshold on the
    information graph) are equally plausible \emph{a priori}.
  \item Proving saturation requires showing that \emph{every} boundary cell
    is independently variable, i.e., that no gauge constraints or
    holographic redundancies reduce the effective number of boundary degrees
    of freedom below $N_{\partial\Omega}$.
  \item The conjecture does \textbf{not} follow from A1--A3 alone.  It
    requires an additional postulate about the relationship between causal
    structure (derived from $q$-gradients) and information capacity.
\end{itemize}

\subsection{Assumption tiering}

\begin{description}
  \item[Tier 0 --- Founding axioms:]
    A1 (information exists), A2 (capacity bounded at scale $\ell$).
    A3 is not strictly required for the inequality --- the min-cut bound
    holds regardless of flow direction.

  \item[Tier 1 --- Theorems (strict):]
    \begin{itemize}
      \item Min-cut entropy bound $S\le N_{\partial\Omega}\ln 2$
        (Theorem~\ref{thm:mincut}): follows from A1~+~A2~+~Ford--Fulkerson.
        This is a \textbf{proven graph-theoretic result}.
      \item Area scaling $S\propto A$ (Eq.~\ref{eq:S4_Narea}): follows from
        $N_{\partial\Omega}=A/\ell^2$ and the continuum limit of the
        information graph.  The scaling is robust; only the coefficient
        depends on microscopic details.
    \end{itemize}

  \item[Tier 2 --- $q$-field input:]
    The $q$-field solution~\eqref{eq:S4_qfield} determines \emph{where}
    the cut lies ($r_s$).  This uses the entropy-variational $q$-field
    equation from the main text, which itself depends on A3 and the
    continuum limit.  The \emph{existence} of a cut does not depend on the
    $q$-field, but its \emph{location} at $r_s$ does.

  \item[Tier 3 --- Model-dependent coefficients and conjectures:]
    \begin{itemize}
      \item Numerical ratio $S_{\rm DGF}/S_{\rm BH}=4\ln 2\,(\ell_P/a)^2$:
        depends on the cutoff scale $a$, which is not independently
        predicted.
      \item Causal Locking Saturation (C1): a conjecture.  Proving or
        refuting C1 is an open problem.
    \end{itemize}
\end{description}

\subsection{Honest summary}

\begin{enumerate}
  \item $S\le N_{\partial\Omega}\propto A$ is a \textbf{theorem}
    (Tier~0+1).  It requires only that information exists (A1) and that
    capacity is bounded per cell (A2).  No $G$, no Einstein equations,
    no horizon thermodynamics.
  \item The \textbf{numerical coefficient} $S_{\rm DGF}/S_{\rm BH}$ is
    \textbf{model-dependent}.  With $a=\ell_P$ it is $\sim 2.8$; exact
    matching with Bekenstein--Hawking fixes $a\approx 1.66\,\ell_P$.
    Both are $O(1)$ and physically reasonable, but DGF does not (yet)
    predict the coefficient from first principles.
  \item The \textbf{$e^{-1/2}$ threshold} for the horizon location
    (Layer~2) is a physically motivated but \emph{non-unique} choice.
    The existence of a cut does not depend on this choice; the location
    $r_s$ does.
  \item \textbf{C1 (Causal Locking Saturation) is a conjecture}, not a
    theorem.  It turns the inequality $S\le N_{\partial\Omega}\ln 2$ into
    the equality $S=N_{\partial\Omega}\ln 2$.  Without C1, the DGF
    entropy bound is an \emph{upper} bound --- it does not yet explain
    \emph{why} black holes saturate it rather than falling somewhere
    below.
  \item The result should be read as: the area scaling of black-hole
    entropy is a necessary consequence of bounded information capacity on
    a graph.  This provides an information-theoretic \emph{origin} for the
    area law, independent of gravitational dynamics.  The exact coefficient
    and saturation mechanism require additional physical input beyond the
    current DGF axioms.
\end{enumerate}

% ============================================================
\begin{acknowledgments}
The author reports no external funding and no competing interests.
\end{acknowledgments}

\begin{thebibliography}{99}

\bibitem{dgf_prd}
H.~Zhongchang, ``Macrocosmic Decoherence as Archive-Decompression
Dynamics,'' DGF v3.1-prd manuscript (2026).

\bibitem{blp2009}
H.-P.~Breuer, E.-M.~Laine, and J.~Piilo,
``Measure for the Degree of Non-Markovian Behavior of Quantum Processes
in Open Systems,'' Phys.\ Rev.\ Lett.\ \textbf{103}, 210401 (2009).

\bibitem{buscemi2014}
F.~Buscemi,
``All Entangled Quantum States Are Nonlocal,''
Phys.\ Rev.\ Lett.\ \textbf{108}, 200401 (2012);
see also F.~Buscemi and N.~Datta,
``General Theory of Information Flow in Quantum Processes,''
Phys.\ Rev.\ A \textbf{93}, 012101 (2016).

\bibitem{ford1956}
L.~R.~Ford~Jr.\ and D.~R.~Fulkerson,
``Maximal Flow Through a Network,''
Can.\ J.\ Math.\ \textbf{8}, 399--404 (1956).

\bibitem{bekenstein1973}
J.~D.~Bekenstein,
``Black Holes and Entropy,''
Phys.\ Rev.\ D \textbf{7}, 2333--2346 (1973).

\bibitem{hawking1975}
S.~W.~Hawking,
``Particle Creation by Black Holes,''
Commun.\ Math.\ Phys.\ \textbf{43}, 199--220 (1975).

\bibitem{greensite1993}
J.~Greensite,
``Dynamical Origin of the Lorentzian Signature of Spacetime,''
arXiv:gr-qc/9210008 (1993).

\end{thebibliography}

\end{document}
